\clearpage

\setlength{\parindent}{2em}

\section{研究计划进度安排及预期目标}

\subsection{进度安排}

本项目按照“需求分析-架构设计-模块开发-系统集成测试”的逻辑主线开展，具体进度安排如下：

\subsubsection{第1-2周：需求分析与系统设计}

\par 完成国内外设计调研与理论建构，深入梳理人机协作范式、上下文感知计算及渐进式披露等相关理论。详实对比现有任务管理工具与生成式AI应用在意图识别精度、多轮状态维持及情感化陪伴方面的典型案例。完成系统需求分析、架构设计和数据库设计，输出《需求规格说明书》和《系统设计文档》。

\subsubsection{第3-4周：后端基础模块开发}

\par 完成后端服务基础架构搭建，包括多渠道登录（短信验证）、用户模块、鉴权模块。实现JSONL流式输出机制，确保AI响应的实时性。完成项目模块基本CRUD和视图联合CRUD功能。

\subsubsection{第5-6周：智能体模块开发}

完成智能体模块核心功能开发，包括子智能体配置、工作流编排、智能路由逻辑、检查点持久化、摘要钩子。实现查询工具集和操作工具集，完成工具调用链可靠性保障机制。

\subsubsection{前端开发与系统集成}

\par 完成移动端前端页面开发，包括首页、任务清单页、AI对话页、日历页、专注页、我的页等核心页面。实现前后端接口联调，完成系统集成。开发桌面小组件，提供贴心提醒功能。

\subsubsection{第9-10周：系统测试与优化}

\par 完成智能体工作流稳定性测试与端到端功能验证，确保多轮对话状态管理的可靠性与工具调用的准确性。完成系统性能基准测试，验证高并发场景下的响应延迟与资源占用情况。优化智能体性能，降低成本、加速响应。

\subsubsection{第11-12周：论文撰写与答辩准备}

\par 完成毕业论文撰写，整理系统设计说明书、测试报告等文档。准备结题答辩材料，包括系统演示、PPT制作和答辩演练。

\subsection{预期目标}

\par （1）系统功能目标：交付一套可部署运行的完整系统，包含移动端应用和后台服务。核心功能包括：任务抽取、AI对话创建任务、子任务管理、专注计时（方糖）、成就系统。

\par （2）技术指标目标：系统响应延迟P95<2秒，工具调用成功率≥95\%，多轮对话状态管理可靠性≥99.5\%，高并发场景下系统稳定运行。

\par （3）文档成果目标：按时高质量交付包含系统设计说明书、可部署运行的完整系统、以及系统测试报告在内的全套成果。完成毕业论文撰写，通过毕业答辩。

\par （4）创新点目标：在移动端碎片化场景下实现低干扰与防挫败设计，建立严谨的“用户意图-任务结构”映射机制，在人机协作式任务构建与长期记忆持久化方面形成系统化突破。

